\chapter{Sources and availability}

Each source is listed with what this research paper takes from it.

\section*{The method and the reading}

\textbf{The method.} The construction this research paper reads with, in its stages of
representation, partition, reference, measure, screening and test oracle, and the rule that a symbol
is reached through equality alone~\cite{src:Quigg-orior-workbook}.

\textbf{Exact ingestion.} The primitive that keeps every digit and imposes no quantum from the reading
end, and the domain-blind readings of which positions are occupied and which carry more than one
value, on which the Pauli count is built. Its use on crystal structures is in the crystallography
research paper~\cite{src:Quigg-crystallography}.

\textbf{The element table.} One transcribed table of atomic number, symbol and the Madelung filling
order, shared with the chemistry research paper.

\textbf{The particle table.} The confirmed Standard Model particles as exact integer quantum numbers,
charge in thirds and spin doubled, with no mass.

\textbf{The equality match rate.} The measure the recurrence is read with, and the coincidence floor a
shuffle reaches.

\section*{Reference data}

\textbf{Element symbols and atomic numbers.} The IUPAC table, the canonical reference that names the
elements 1 to 118. It is standard and carries no single locator; the element table transcribes it and
states so.

\textbf{The Madelung filling order and Hund's rule.} Standard atomic structure. The order fills 1s
through 7p and its capacities close on 118, which the element table checks and does not assert.

\section*{The measured data}

\textbf{NIST Atomic Spectra Database}~\cite{src:NIST-Atomic-Spectra-Database}. The measured
ground-state configurations of neutral atoms through element 108, which hold the aufbau exceptions the
ideal filling does not. The comparison against them is not reported here.

\textbf{The predicted superheavy configurations.} Elements 109 to 118 have no measured neutral atom;
their calculated configurations are compiled
in~\cite{src:Wikipedia-Electron-configurations-of-the-elements-data-page}, marked predicted and kept
apart from the measured ones.

\textbf{The Particle Data Group}~\cite{src:Particle-Data-Group}. The confirmed particle content, three
generations and the Higgs boson, and the measured masses, which the particle table records and this
reading sets aside.

\textbf{CODATA recommended values}~\cite{src:Mohr-et-al-2025}. The Rydberg constant, the Bohr radius
and the fine-structure constant, the measured hubs that set the absolute scale the exact rationals
stand on. The trajectory chapter's absolute comparison needs them and is not run; no value is quoted.

\section*{Availability}

The representation, the measures and the runs that print every figure in this research paper are in
the orior repository, \texttt{https://github.com/dstroy0/orior}.
